%Root main.tex
%---------------------------------------------------------------------------
%付録A　付録
%---------------------------------------------------------------------------

%\lstnewenvironment{mylisting}[1][]
%    {\lstset{
%        frame=single,
%        numbersep=10pt,
%        tabsize=2,
%        #1
%    }
%}{}
\begin{mylisting}[language=c,caption=電流方向検出モジュール,label=hurokuB_denryu]

#include <Stdlib.h>
#include <String.h>
#include <math.h>
#include <Psim.h>

// PLACE GLOBAL VARIABLES OR USER FUNCTIONS HERE...

double iL;
int K;

/////////////////////////////////////////////////////////////////////
// FUNCTION: SimulationStep
//   This function runs at every time step.
//double t: (read only) time
//double delt: (read only) time step as in Simulation control
//double *in: (read only) zero based array of input values. in[0] is the first node, in[1] second input...
//double *out: (write only) zero based array of output values. out[0] is the first node, out[1] second output...
//int *pnError: (write only)  assign  *pnError = 1;  if there is an error and set the error message in szErrorMsg
//    strcpy(szErrorMsg, "Error message here..."); 
// DO NOT CHANGE THE NAME OR PARAMETERS OF THIS FUNCTION
void SimulationStep(
		double t, double delt, double *in, double *out,
		 int *pnError, char * szErrorMsg,
		 void ** reserved_UserData, int reserved_ThreadIndex, void * reserved_AppPtr)
{
// ENTER YOUR CODE HERE...
//in
iL = in[0];

	if(iL >= 0){
	K = 0;
	}else{
	K = 1;
	}

//out
out[0] =K;
}


/////////////////////////////////////////////////////////////////////
// FUNCTION: SimulationBegin
//   Initialization function. This function runs once at the beginning of simulation
//   For parameter sweep or AC sweep simulation, this function runs at the beginning of each simulation cycle.
//   Use this function to initialize static or global variables.
//const char *szId: (read only) Name of the C-block 
//int nInputCount: (read only) Number of input nodes
//int nOutputCount: (read only) Number of output nodes
//int nParameterCount: (read only) Number of parameters is always zero for C-Blocks.  Ignore nParameterCount and pszParameters
//int *pnError: (write only)  assign  *pnError = 1;  if there is an error and set the error message in szErrorMsg
//    strcpy(szErrorMsg, "Error message here..."); 
// DO NOT CHANGE THE NAME OR PARAMETERS OF THIS FUNCTION
void SimulationBegin(
		const char *szId, int nInputCount, int nOutputCount,
		 int nParameterCount, const char ** pszParameters,
		 int *pnError, char * szErrorMsg,
		 void ** reserved_UserData, int reserved_ThreadIndex, void * reserved_AppPtr)
{
// ENTER INITIALIZATION CODE HERE...




}


/////////////////////////////////////////////////////////////////////
// FUNCTION: SimulationEnd
//   Termination function. This function runs once at the end of simulation
//   For parameter sweep or AC sweep simulation, this function runs at the end of each simulation cycle.
//   Use this function to de-allocate any allocated memory or to save the result of simulation in an alternate file.
// Ignore all parameters for C-block 
// DO NOT CHANGE THE NAME OR PARAMETERS OF THIS FUNCTION
void SimulationEnd(const char *szId, void ** reserved_UserData, int reserved_ThreadIndex, void * reserved_AppPtr)
{




}





\end{mylisting}

\begin{mylisting}[language=c,caption=1番レグ用比較モジュール,label=hurokuB_hikakuA]

#include <Stdlib.h>
#include <String.h>
#include <math.h>
#include <Psim.h>

// PLACE GLOBAL VARIABLES OR USER FUNCTIONS HERE...
double in1,in2,in3,in4;
int buf1,buf2,buf3;
int out1;

/////////////////////////////////////////////////////////////////////
// FUNCTION: SimulationStep
//   This function runs at every time step.
//double t: (read only) time
//double delt: (read only) time step as in Simulation control
//double *in: (read only) zero based array of input values. in[0] is the first node, in[1] second input...
//double *out: (write only) zero based array of output values. out[0] is the first node, out[1] second output...
//int *pnError: (write only)  assign  *pnError = 1;  if there is an error and set the error message in szErrorMsg
//    strcpy(szErrorMsg, "Error message here..."); 
// DO NOT CHANGE THE NAME OR PARAMETERS OF THIS FUNCTION
void SimulationStep(
		double t, double delt, double *in, double *out,
		 int *pnError, char * szErrorMsg,
		 void ** reserved_UserData, int reserved_ThreadIndex, void * reserved_AppPtr)
{
// ENTER YOUR CODE HERE...
//in

in1 = in[0];
in2 = in[1];
in3 = in[2];
in4 = in[3];

	if(in1 > in2){
	buf1 = 1;
	}else if(in1 == in2){
	buf1 = 1;
	}else{
	buf1 = 0;
	}

	if(in1 >= in3){
	buf2 = 1;
	}else if(in1 == in3){
	buf1 = 1;
	}else{
	buf2 = 0;
	}
	if(in1 >= in4){
	buf3 = 1;
	}else if(in1 == in3){
	buf1 = 1;
	}else{
	buf3 = 0;
	}
	
	//out1 = buf1 + buf2 + buf3;

//out
out[0] =buf1;
out[1] =buf2;
out[2] =buf3;
}


/////////////////////////////////////////////////////////////////////
// FUNCTION: SimulationBegin
//   Initialization function. This function runs once at the beginning of simulation
//   For parameter sweep or AC sweep simulation, this function runs at the beginning of each simulation cycle.
//   Use this function to initialize static or global variables.
//const char *szId: (read only) Name of the C-block 
//int nInputCount: (read only) Number of input nodes
//int nOutputCount: (read only) Number of output nodes
//int nParameterCount: (read only) Number of parameters is always zero for C-Blocks.  Ignore nParameterCount and pszParameters
//int *pnError: (write only)  assign  *pnError = 1;  if there is an error and set the error message in szErrorMsg
//    strcpy(szErrorMsg, "Error message here..."); 
// DO NOT CHANGE THE NAME OR PARAMETERS OF THIS FUNCTION
void SimulationBegin(
		const char *szId, int nInputCount, int nOutputCount,
		 int nParameterCount, const char ** pszParameters,
		 int *pnError, char * szErrorMsg,
		 void ** reserved_UserData, int reserved_ThreadIndex, void * reserved_AppPtr)
{
// ENTER INITIALIZATION CODE HERE...




}


/////////////////////////////////////////////////////////////////////
// FUNCTION: SimulationEnd
//   Termination function. This function runs once at the end of simulation
//   For parameter sweep or AC sweep simulation, this function runs at the end of each simulation cycle.
//   Use this function to de-allocate any allocated memory or to save the result of simulation in an alternate file.
// Ignore all parameters for C-block 
// DO NOT CHANGE THE NAME OR PARAMETERS OF THIS FUNCTION
void SimulationEnd(const char *szId, void ** reserved_UserData, int reserved_ThreadIndex, void * reserved_AppPtr)
{




}





\end{mylisting}

\begin{mylisting}[language=c,caption=2番レグ用比較モジュール,label=hurokuB_hikakuB]

#include <Stdlib.h>
#include <String.h>
#include <math.h>
#include <Psim.h>

// PLACE GLOBAL VARIABLES OR USER FUNCTIONS HERE...
double in1,in2,in3,in4;
int buf1,buf2,buf3;
int out1;

/////////////////////////////////////////////////////////////////////
// FUNCTION: SimulationStep
//   This function runs at every time step.
//double t: (read only) time
//double delt: (read only) time step as in Simulation control
//double *in: (read only) zero based array of input values. in[0] is the first node, in[1] second input...
//double *out: (write only) zero based array of output values. out[0] is the first node, out[1] second output...
//int *pnError: (write only)  assign  *pnError = 1;  if there is an error and set the error message in szErrorMsg
//    strcpy(szErrorMsg, "Error message here..."); 
// DO NOT CHANGE THE NAME OR PARAMETERS OF THIS FUNCTION
void SimulationStep(
		double t, double delt, double *in, double *out,
		 int *pnError, char * szErrorMsg,
		 void ** reserved_UserData, int reserved_ThreadIndex, void * reserved_AppPtr)
{
// ENTER YOUR CODE HERE...
//in

in1 = in[0];
in2 = in[1];
in3 = in[2];
in4 = in[3];

	if(in1 > in2){
	buf1 = 1;
	}else if(in1 == in2){
	buf1 = 0;
	}else{
	buf1 = 0;
	}

	if(in1 >= in3){
	buf2 = 1;
	}else if(in1 == in3){
	buf1 = 1;
	}else{
	buf2 = 0;
	}
	if(in1 >= in4){
	buf3 = 1;
	}else if(in1 == in3){
	buf1 = 1;
	}else{
	buf3 = 0;
	}
	
	//out1 = buf1 + buf2 + buf3;

//out
out[0] =buf1;
out[1] =buf2;
out[2] =buf3;
}


/////////////////////////////////////////////////////////////////////
// FUNCTION: SimulationBegin
//   Initialization function. This function runs once at the beginning of simulation
//   For parameter sweep or AC sweep simulation, this function runs at the beginning of each simulation cycle.
//   Use this function to initialize static or global variables.
//const char *szId: (read only) Name of the C-block 
//int nInputCount: (read only) Number of input nodes
//int nOutputCount: (read only) Number of output nodes
//int nParameterCount: (read only) Number of parameters is always zero for C-Blocks.  Ignore nParameterCount and pszParameters
//int *pnError: (write only)  assign  *pnError = 1;  if there is an error and set the error message in szErrorMsg
//    strcpy(szErrorMsg, "Error message here..."); 
// DO NOT CHANGE THE NAME OR PARAMETERS OF THIS FUNCTION
void SimulationBegin(
		const char *szId, int nInputCount, int nOutputCount,
		 int nParameterCount, const char ** pszParameters,
		 int *pnError, char * szErrorMsg,
		 void ** reserved_UserData, int reserved_ThreadIndex, void * reserved_AppPtr)
{
// ENTER INITIALIZATION CODE HERE...




}


/////////////////////////////////////////////////////////////////////
// FUNCTION: SimulationEnd
//   Termination function. This function runs once at the end of simulation
//   For parameter sweep or AC sweep simulation, this function runs at the end of each simulation cycle.
//   Use this function to de-allocate any allocated memory or to save the result of simulation in an alternate file.
// Ignore all parameters for C-block 
// DO NOT CHANGE THE NAME OR PARAMETERS OF THIS FUNCTION
void SimulationEnd(const char *szId, void ** reserved_UserData, int reserved_ThreadIndex, void * reserved_AppPtr)
{




}





\end{mylisting}

\begin{mylisting}[language=c,caption=3番レグ用比較モジュール,label=hurokuB_hikakuC]

#include <Stdlib.h>
#include <String.h>
#include <math.h>
#include <Psim.h>

// PLACE GLOBAL VARIABLES OR USER FUNCTIONS HERE...
double in1,in2,in3,in4;
int buf1,buf2,buf3;
int out1;

/////////////////////////////////////////////////////////////////////
// FUNCTION: SimulationStep
//   This function runs at every time step.
//double t: (read only) time
//double delt: (read only) time step as in Simulation control
//double *in: (read only) zero based array of input values. in[0] is the first node, in[1] second input...
//double *out: (write only) zero based array of output values. out[0] is the first node, out[1] second output...
//int *pnError: (write only)  assign  *pnError = 1;  if there is an error and set the error message in szErrorMsg
//    strcpy(szErrorMsg, "Error message here..."); 
// DO NOT CHANGE THE NAME OR PARAMETERS OF THIS FUNCTION
void SimulationStep(
		double t, double delt, double *in, double *out,
		 int *pnError, char * szErrorMsg,
		 void ** reserved_UserData, int reserved_ThreadIndex, void * reserved_AppPtr)
{
// ENTER YOUR CODE HERE...
//in

in1 = in[0];
in2 = in[1];
in3 = in[2];
in4 = in[3];

	if(in1 > in2){
	buf1 = 1;
	}else if(in1 == in2){
	buf1 = 0;
	}else{
	buf1 = 0;
	}

	if(in1 >= in3){
	buf2 = 1;
	}else if(in1 == in3){
	buf1 = 0;
	}else{
	buf2 = 0;
	}
	if(in1 >= in4){
	buf3 = 1;
	}else if(in1 == in3){
	buf1 = 1;
	}else{
	buf3 = 0;
	}
	
	//out1 = buf1 + buf2 + buf3;

//out
out[0] =buf1;
out[1] =buf2;
out[2] =buf3;
}


/////////////////////////////////////////////////////////////////////
// FUNCTION: SimulationBegin
//   Initialization function. This function runs once at the beginning of simulation
//   For parameter sweep or AC sweep simulation, this function runs at the beginning of each simulation cycle.
//   Use this function to initialize static or global variables.
//const char *szId: (read only) Name of the C-block 
//int nInputCount: (read only) Number of input nodes
//int nOutputCount: (read only) Number of output nodes
//int nParameterCount: (read only) Number of parameters is always zero for C-Blocks.  Ignore nParameterCount and pszParameters
//int *pnError: (write only)  assign  *pnError = 1;  if there is an error and set the error message in szErrorMsg
//    strcpy(szErrorMsg, "Error message here..."); 
// DO NOT CHANGE THE NAME OR PARAMETERS OF THIS FUNCTION
void SimulationBegin(
		const char *szId, int nInputCount, int nOutputCount,
		 int nParameterCount, const char ** pszParameters,
		 int *pnError, char * szErrorMsg,
		 void ** reserved_UserData, int reserved_ThreadIndex, void * reserved_AppPtr)
{
// ENTER INITIALIZATION CODE HERE...




}


/////////////////////////////////////////////////////////////////////
// FUNCTION: SimulationEnd
//   Termination function. This function runs once at the end of simulation
//   For parameter sweep or AC sweep simulation, this function runs at the end of each simulation cycle.
//   Use this function to de-allocate any allocated memory or to save the result of simulation in an alternate file.
// Ignore all parameters for C-block 
// DO NOT CHANGE THE NAME OR PARAMETERS OF THIS FUNCTION
void SimulationEnd(const char *szId, void ** reserved_UserData, int reserved_ThreadIndex, void * reserved_AppPtr)
{




}





\end{mylisting}

\begin{mylisting}[language=c,caption=4番レグ用比較モジュール,label=hurokuB_hikakuD]

#include <Stdlib.h>
#include <String.h>
#include <math.h>
#include <Psim.h>

// PLACE GLOBAL VARIABLES OR USER FUNCTIONS HERE...
double in1,in2,in3,in4;
int buf1,buf2,buf3;
int out1;

/////////////////////////////////////////////////////////////////////
// FUNCTION: SimulationStep
//   This function runs at every time step.
//double t: (read only) time
//double delt: (read only) time step as in Simulation control
//double *in: (read only) zero based array of input values. in[0] is the first node, in[1] second input...
//double *out: (write only) zero based array of output values. out[0] is the first node, out[1] second output...
//int *pnError: (write only)  assign  *pnError = 1;  if there is an error and set the error message in szErrorMsg
//    strcpy(szErrorMsg, "Error message here..."); 
// DO NOT CHANGE THE NAME OR PARAMETERS OF THIS FUNCTION
void SimulationStep(
		double t, double delt, double *in, double *out,
		 int *pnError, char * szErrorMsg,
		 void ** reserved_UserData, int reserved_ThreadIndex, void * reserved_AppPtr)
{
// ENTER YOUR CODE HERE...
//in

in1 = in[0];
in2 = in[1];
in3 = in[2];
in4 = in[3];

	if(in1 > in2){
	buf1 = 1;
	}else if(in1 == in2){
	buf1 = 0;
	}else{
	buf1 = 0;
	}

	if(in1 >= in3){
	buf2 = 1;
	}else if(in1 == in3){
	buf1 = 0;
	}else{
	buf2 = 0;
	}
	if(in1 >= in4){
	buf3 = 1;
	}else if(in1 == in3){
	buf1 = 0;
	}else{
	buf3 = 0;
	}
	
	//out1 = buf1 + buf2 + buf3;

//out
out[0] =buf1;
out[1] =buf2;
out[2] =buf3;
}


/////////////////////////////////////////////////////////////////////
// FUNCTION: SimulationBegin
//   Initialization function. This function runs once at the beginning of simulation
//   For parameter sweep or AC sweep simulation, this function runs at the beginning of each simulation cycle.
//   Use this function to initialize static or global variables.
//const char *szId: (read only) Name of the C-block 
//int nInputCount: (read only) Number of input nodes
//int nOutputCount: (read only) Number of output nodes
//int nParameterCount: (read only) Number of parameters is always zero for C-Blocks.  Ignore nParameterCount and pszParameters
//int *pnError: (write only)  assign  *pnError = 1;  if there is an error and set the error message in szErrorMsg
//    strcpy(szErrorMsg, "Error message here..."); 
// DO NOT CHANGE THE NAME OR PARAMETERS OF THIS FUNCTION
void SimulationBegin(
		const char *szId, int nInputCount, int nOutputCount,
		 int nParameterCount, const char ** pszParameters,
		 int *pnError, char * szErrorMsg,
		 void ** reserved_UserData, int reserved_ThreadIndex, void * reserved_AppPtr)
{
// ENTER INITIALIZATION CODE HERE...




}


/////////////////////////////////////////////////////////////////////
// FUNCTION: SimulationEnd
//   Termination function. This function runs once at the end of simulation
//   For parameter sweep or AC sweep simulation, this function runs at the end of each simulation cycle.
//   Use this function to de-allocate any allocated memory or to save the result of simulation in an alternate file.
// Ignore all parameters for C-block 
// DO NOT CHANGE THE NAME OR PARAMETERS OF THIS FUNCTION
void SimulationEnd(const char *szId, void ** reserved_UserData, int reserved_ThreadIndex, void * reserved_AppPtr)
{




}





\end{mylisting}

\begin{mylisting}[language=c,caption=各スイッチングモジュール,label=hurokuB_switxhhikaku]

#include <Stdlib.h>
#include <String.h>
#include <math.h>
#include <Psim.h>

// PLACE GLOBAL VARIABLES OR USER FUNCTIONS HERE...
double sw;
int in1,in2,in3,in4;
int out1,out2,out3,out4;

/////////////////////////////////////////////////////////////////////
// FUNCTION: SimulationStep
//   This function runs at every time step.
//double t: (read only) time
//double delt: (read only) time step as in Simulation control
//double *in: (read only) zero based array of input values. in[0] is the first node, in[1] second input...
//double *out: (write only) zero based array of output values. out[0] is the first node, out[1] second output...
//int *pnError: (write only)  assign  *pnError = 1;  if there is an error and set the error message in szErrorMsg
//    strcpy(szErrorMsg, "Error message here..."); 
// DO NOT CHANGE THE NAME OR PARAMETERS OF THIS FUNCTION
void SimulationStep(
		double t, double delt, double *in, double *out,
		 int *pnError, char * szErrorMsg,
		 void ** reserved_UserData, int reserved_ThreadIndex, void * reserved_AppPtr)
{
// ENTER YOUR CODE HERE...
//in
sw = in[0];
in1 = in[1];
in2 = in[2];
in3 = in[3];
in4 = in[4];

	if(in1 >= sw){
	out1 = 1;
	}else{
	out1 = 0;
	}

	if(in2 >= sw){
	out2 = 1;
	}else{
	out2 = 0;
	}

	if(in3 >= sw){
	out3 = 1;
	}else{
	out3 = 0;
	}

	if(in4 >= sw){
	out4 = 1;
	}else{
	out4 = 0;
	}

//out
out[0] =out1;
out[1] =out2;
out[2] =out3;
out[3] =out4;
}


/////////////////////////////////////////////////////////////////////
// FUNCTION: SimulationBegin
//   Initialization function. This function runs once at the beginning of simulation
//   For parameter sweep or AC sweep simulation, this function runs at the beginning of each simulation cycle.
//   Use this function to initialize static or global variables.
//const char *szId: (read only) Name of the C-block 
//int nInputCount: (read only) Number of input nodes
//int nOutputCount: (read only) Number of output nodes
//int nParameterCount: (read only) Number of parameters is always zero for C-Blocks.  Ignore nParameterCount and pszParameters
//int *pnError: (write only)  assign  *pnError = 1;  if there is an error and set the error message in szErrorMsg
//    strcpy(szErrorMsg, "Error message here..."); 
// DO NOT CHANGE THE NAME OR PARAMETERS OF THIS FUNCTION
void SimulationBegin(
		const char *szId, int nInputCount, int nOutputCount,
		 int nParameterCount, const char ** pszParameters,
		 int *pnError, char * szErrorMsg,
		 void ** reserved_UserData, int reserved_ThreadIndex, void * reserved_AppPtr)
{
// ENTER INITIALIZATION CODE HERE...




}


/////////////////////////////////////////////////////////////////////
// FUNCTION: SimulationEnd
//   Termination function. This function runs once at the end of simulation
//   For parameter sweep or AC sweep simulation, this function runs at the end of each simulation cycle.
//   Use this function to de-allocate any allocated memory or to save the result of simulation in an alternate file.
// Ignore all parameters for C-block 
// DO NOT CHANGE THE NAME OR PARAMETERS OF THIS FUNCTION
void SimulationEnd(const char *szId, void ** reserved_UserData, int reserved_ThreadIndex, void * reserved_AppPtr)
{




}





\end{mylisting}

